% TOP_PROBLEM: {"id": 4600017, "title": "Furstenberg times-p-times-q problem", "category_id": 11, "category": {"id": 11, "name": "dynamical_systems", "display_name": "Dynamical Systems", "description": "Problems about long-term behavior of deterministic systems, Hamiltonian dynamics, and periodic orbits.", "slug": "dynamical-systems", "order_index": 11, "created_at": "2026-07-31T15:26:25.671Z"}, "status": "partially_solved", "problem_number": "AMR-045-0017", "set_id": 13}
% FIRST_POSTED: 2026-10-01
% ENTRY_KIND: statement_only
% UnsolvedMath ID 4600017; source catalogue code AMR-045-0017
% !TeX program = lualatex
% Definition and sources only; this file is not an LLM solution attempt.
\documentclass[11pt,a4paper]{article}
\usepackage[margin=25mm]{geometry}
\usepackage{fontspec}
\setmainfont{lmroman10-regular.otf}[BoldFont=lmroman10-bold.otf,
 ItalicFont=lmroman10-italic.otf,BoldItalicFont=lmroman10-bolditalic.otf]
\usepackage{amsmath,amssymb,mathtools}
\usepackage{unicode-math}
\setmathfont{latinmodern-math.otf}
\AtBeginDocument{\let\setminus\smallsetminus}
\usepackage{xurl,hyperref}
\hypersetup{colorlinks=true,urlcolor=blue}
\setcounter{secnumdepth}{0}
\setlength{\parindent}{0pt}
\setlength{\parskip}{5pt}
\setlength{\emergencystretch}{3em}
\begin{document}
\section*{Furstenberg times-p-times-q problem}\label{4600017}
{\small UnsolvedMath ID 4600017 · AMR-045-0017 · Dynamical Systems · AMR Open Problem Lists}
\subsection{Definitions and mathematical statement}
Write $\mathbb T=\mathbb R/\mathbb Z$ for the circle with its Borel
$\sigma$-algebra, and let $m$ be normalized Lebesgue (Haar) measure.
Fix integers $p,q>1$ which are multiplicatively independent, meaning
$p^r\ne q^s$ for all positive integers $r,s$. Define the continuous
endomorphisms $M_p(x)=px\pmod1$ and $M_q(x)=qx\pmod1$.

Suppose $\mu$ is a Borel probability on $\mathbb T$ satisfying
\[
 \mu(M_p^{-1}E)=\mu(E)=\mu(M_q^{-1}E)
 \qquad\text{for every Borel }E\subset\mathbb T,
\]
and $\mu(\{x\})=0$ for every $x\in\mathbb T$.
Must $\mu=m$? Equivalently, can there be a nonatomic probability other
than Haar measure invariant under both multiplications?

Invariance is expressed by inverse images because $M_p$ and $M_q$ are
not one-to-one maps of the circle. The problem puts no absolute continuity
or full-support assumption on $\mu$ and no lower bound on its entropy.
It also does not assume ergodicity for either individual map. The
nonatomic hypothesis excludes measures supported on finite rational
orbits and all probability measures assigning positive mass to a point.
Multiplicative independence excludes the immediate situation in which
both maps are powers of the same multiplication map. Both the integers
and the measure are universally quantified in the uniqueness formulation;
a counterexample for any one admissible pair would answer the existence
question affirmatively.
\subsection{Short English statement}
Is Lebesgue measure the only nonatomic probability on the circle invariant under multiplication by two multiplicatively independent integers?
\subsection{Sources}
\begin{itemize}
\item[S1] ULAM.AI and the UnsolvedMath Contributors, supplied problems.json, record 4600017 (AMR-045-0017), AMR Open Problem Lists. Catalogue identity and source transcription; CC BY 4.0.
\url{https://huggingface.co/datasets/ulamai/UnsolvedMath}.
\item[S2] Boyle - Open Problems in Symbolic Dynamics (2008), Problem/Question/Conjecture 14.1. Original source indexed by AMR-045-0017.
\url{https://www.math.umd.edu/~mboyle/papers/openfinalsub3nov2007.pdf}.
\end{itemize}
\end{document}
